\documentstyle[%


righttag%


,11pt%


,amssymb%


,russian%               remove in Eng.


,izvru%                 Set izven% in Eng.


]{amsart}





%


\newcommand{\Mod}{\operatorname{mod}}


\newcommand{\tts}[1]{}


\renewcommand{\baselinestretch}{0.96}





\theoremstyle{plain}


\theorembodyfont{\it}


\newtheorem{thmnonum}{\hskip\parindent Теорема}


\newtheorem{propnonum}{\hskip\parindent Предложение}


\renewcommand{\thethmnonum}{}


\renewcommand{\thepropnonum}{}





\theoremstyle{definition}


\theorembodyfont{\rm}


\newtheorem{defnnonum}{\hskip\parindent Определение}


\newtheorem{notenonum}{\hskip\parindent Замечание}


\renewcommand{\thedefnnonum}{}


\renewcommand{\thenotenonum}{}





%\hoffset=-15mm%                  For Eng.


\setcounter{page}{47}


\god{2001}


\nomer{3~(466)} %                        Remove in Eng.


%\nomer{Vol.\,45, No.\,3}%               Set in Eng.


%\str{\pageref{begin}--\pageref{end}}%  Set in Eng.


\udk{514.763}





\author{\it К.В.\,Семенов}


% NB:   ^^ remove in Eng.


\title{О геометрии псевдопотенциалов эволюционного


уравнения третьего порядка}





\date{}


%\pagestyle{myheadings}%         Set in Eng.


\pagestyle{plain} %              Remove in Eng.


\begin{document}


%\thispagestyle{plain}%          Set in Eng.


\maketitle


\label{begin}





% enum. of eqs. 





\section*{Введение}





В данной работе рассматривается геометрия эволюционного 


уравнения третьего порядка


и вопрос о существовании псевдопотенциалов. 


При рассмотрении вопроса о существовании связности, 


определяющей представление нулевой


кривизны для уравнения, применяется метод, предложенный в [1]


для изучения эволюционных уравнений второго порядка.


Полученные результаты обобщают результаты статьи [1]


на случай уравнений более


высокого порядка.





\section{Геометрия эволюционного дифференциального уравнения третьего порядка.


Фундаментальный объект}





Объектом исследования в данной работе являются 


дифференциальные уравнения с частными производными


третьего порядка


$$


u_t=f(t,x^1,\dotsc,x^n,u,u_j,u_{jk},u_{jkl})\quad (j,k,l=1,\dotsc,n).


$$


Здесь $u$ --- неизвестная функция, $u_t$ --- ее производная 


по переменной $t$, имеющей


смысл времени, $u_i$ --- частные производные по пространственным 


переменным; функция $f$


не содержит $u_t$.


Уравнения такого вида носят название эволюционных. 


Таковым является известное уравнение


Кортевега--де~Фриза.





Будем полагать, что $(t,x^1,\dotsc ,x^n,u,u_j,u_{jk},u_{jkl})$ ---


адаптированные локальные координаты $(n+2)$-мерного расслоения $E$ с


(расслоенной) $(n+1)$-мерной базой $M$;


координаты в $M$ --- переменные $t$, $x^1,\dotsc,x^n$.


Допустимые преобразования координат задаются формулами


\begin{equation}


\begin{cases}


{\wt t}={{\varphi}^0(t)}, & \ \\


{\wt x}^i=\varphi^i (t,x^1, \dotsc, x^n), \\


{\wt u}={\varphi^{n+1}(t,x^1, \dotsc, x^n, u)}.


\end{cases}


\quad i=1,\dotsc, n, 


\tag{1а}


\end{equation}


Bыражая $t$, $x^1,\dotsc,x^n$, $u$ через


$\wt{t}$, $\wt{x}^1,\dotsc,\wt{x}^n$, $\wt{u}$, получим аналогичные формулы


\begin{equation}


\begin{cases}


t=\psi^0({\wt t}), \\


x^i={\psi}^i({\wt t}, {\wt x}^1, \dotsc, {\wt x}^n),\\


u={\psi}^{n+1} ({\wt t}, {\wt x}^1, \dotsc, {\wt x}^n, \wt u).


\end{cases}


\quad i=1,\dotsc,n,


\tag{1б}


\end{equation}


\addtocounter{equation}{1}





Переходя к расслоению $J^3E$ голономных струй третьего порядка


расслоения $E$, наше уравнение запишем в виде


\begin{equation}


{\lambda_0}-f(t,x^1,\dotsc,x^n,u,{\lambda_j},


{\lambda_{jk}},{\lambda_{jkl}})=0.


\end{equation}


Закон преобразования


$\lambda_{\Hat{i}}$, $\lambda_{\Hat{i}\Hat{k}}$,


$\lambda_{\Hat{i}\Hat{k}\Hat{l}}$ такой же, как


у соответствующих частных производных


$u_{\Hat{i}}$, $u_{\Hat{i}\Hat{k}}$,


$u_{\Hat{i}\Hat{k}\Hat{l}}$


$(\wh{i},\wh{k},\wh{l}=0,\dotsc,n)$


сечения расслоенного многообразия, заданного соотношениями


$u=u(t,x^1,\dotsc ,x^n)$. Допустимые преобразования для координаты $\lambda_0$


записываются следующим образом:


$$


\lambda_0 =


\frac{\partial\psi^{n+1}}{ \partial \wt u}


\frac{d\varphi^0}{d t}\wt\lambda_0+


\frac{\partial\psi^{n+1}}{ \partial \wt u}


\frac{d\varphi^j}{d t}\wt\lambda_j+


\frac{\partial\psi^{n+1}}{ \partial \wt t}


\frac{d\varphi^0}{d t}+


\frac{\partial\psi^{n+1}}{ \partial \wt x^j}


\frac{d\varphi^j}{d t}.


$$


Нетрудно выписать и допустимые преобразования


координат $\lambda_i$, $\lambda_{ij}$, $\lambda_{ijk}.$


Анализируя соотношения (1б) и выражения для допустимых преобразований


координат


$\lambda_i$, $\lambda_{ij}$, $\lambda_{ijk},$ можно заключить, что


переменные


$t$, $x^i$, $u$, $\lambda_{\Hat{j}}$, $\lambda_{ij}$, $\lambda_{ijk}$ могут


одновременно рассматриваться и как часть локальных координат


многообразия $J^3E$, и как


вся совокупность локальных координат некоторого


многообразия $\ov{J}^3E$, которое


является фактормногообразием многообразия $J^3E$.





Ранее [1] уже отмечалось, что уравнение (2),


вообще говоря, теряет специальный вид


при произвольных преобразованиях координат. Однако возможно указать некоторые


дополнительные условия, при которых вид уравнения не меняется.


В результате допустимых преобразований уравнение (2) принимает вид


$$


\frac{\partial\psi^{n+1}}{\partial\wt{u}}


\frac{d\varphi^0}{dt}\wt{\lambda_0}+\dotsb-f


\bigg|_{t=\psi^0(\Tilde{t}), \


x^i=\psi^i(\Tilde{x}^1,\dotsc,\Tilde{x}^n),\dotsc}


=0.


$$


Первоначальный вид уравнения не изменится, если допустимые преобразования


удовлетворяют условию


$$


\frac{\partial\psi^{n+1}}{\partial\wt{u}}


\frac{d\varphi^0}{dt}=1,


$$


что имеет место, например, в случае


\begin{equation}


\begin{cases}


t=\wt{t}+c, \\


x^i=\varphi^i(t,\wt{x}^1,\dotsc,\wt{x}^n), \\


u=\wt{u}+\psi(t,\wt{x}^1,\dotsc,\wt{x}^n).


\end{cases}


\end{equation}





Можно утверждать, что задание эволюционного уравнения (2) 


эквивалентно заданию на


многообразии $\ov{J}^3E$ функции вида


$$


\lambda_0-f(t,{x}^i,u,\lambda_j,\lambda_{kl},\lambda_{jkl})


$$


при условии дополнительных ограничений (3) на преобразования координат


на $E$.





Пусть


${\omega }^{\Hat i}$, ${\omega }^{n+1}$, ${\omega }^{n+1}_{n+1}$,


${\omega }^0_0$,


${\omega }^i_{\Hat j}$, ${\omega }^{n+1}_{j,{n+1}}$,


${\omega }^{n+1}_{\Hat j}$,


${\omega }^0_{00}$,


${\omega }^i_{{\Hat j}{\Hat k}}$, ${\omega }^{n+1}_{{n+1},{n+1}}, \dotsc$


и т.\,д. ---


cимметричные по нижним индексам структурные формы расслоений


голономных реперов многообразия $E$.


При этом


${\omega }^{\Hat i}$, ${\omega }^{n+1}$,


${\omega }^{n+1}_{{\Hat i}_1\dotsc{\Hat i}_k}$


($k=1,2,\dotsc,r$)


--- главные формы в расслоении голономных $r$-струй $J^rE$.


Например,


$$


\omega^{\Hat{i}}=dx^{\Hat{i}}\ \; (\omega^0=dt),\ \;


\omega^{n+1}=du-\lambda_{\Hat{i}}dx^{\Hat{i}}, \ \;


\omega^{n+1}_{\Hat{i}_1\dotsc\Hat{i}_k}=


d\lambda_{\Hat{i}_1\dotsc\Hat{i}_k}-


\lambda_{\Hat{i}_1\dotsc\Hat{i}_k\Hat{j}}dx^{\Hat{j}}


$$


--- контактные формы.





Структурные уравнения, которым они удовлетворяют, суть


\begin{equation}


\begin{cases}


d\omega^0=\omega^0\wedge\omega^0_0, \\


d\omega^i=\omega^0\wedge\omega^i_0+\omega^j\wedge\omega^i_j,\\


d\omega^{n+1}=\omega^0\wedge\omega^{n+1}_0+\omega^j\wedge


\omega^{n+1}_j+\omega^{n+1}\wedge\omega^{n+1}_{n+1}


\end{cases}


\end{equation}


и уравнения, получающиеся продолжением (4).





Запишем дифференциал функции


$$


{\lambda_0}-f(t,x^1,\dotsc,x^n,u,{\lambda_j},


{\lambda_{jk}},{\lambda_{jkl}})


$$


(левой части уравнения (2)) в виде


$$


d(\lambda_0-f)=\omega_0^{n+1}-F_0\omega^0-


F_i\omega^i-F\omega^{n+1}-F^j\omega^{n+1}_j-F^{jk}\omega_{jk}^{n+1}


-F^{jkl}\omega_{jkl}^{n+1}.


$$





Продолжая это соотношение, получим дифференциальные уравнения, 


которым удовлетворяют коэффициенты $F^{jkl}$, $F^{jk}$, $F^j$, $F$,


$F_i$, $F_0$:


\begin{align*}


dF^{jkl}+F^{mkl}{\omega }_m^j+F^{jml}{\omega }_m^k+


F^{jkm}{\omega }^l_m-F^{jkl}{\omega }^{n+1}_{n+1}&=0, \\


%


dF^{jk}+F^{mk}{\omega }^j_m+F^{jm}{\omega }_m^k-


F^{jk}{\omega }^{n+1}_{n+1}+\frac32 F^{uv(j}{\omega }^{k)}_{uv}-


3F^{jkm}{\omega }^{n+1}_{m,n+1}&=0, \\


%


dF^i+F^m{\omega }^i_m -F^i{\omega }^{n+1}_{n+1}-


2F^{im}{\omega }^{n+1}_{m,n+1}


+F^{uv}{\omega }^i_{uv}+F^{uvw}{\omega }^i_{uvw} -


3F^{iuv}{\omega }^{n+1}_{uv,n+1}&=0, \\


%


dF-F{\omega }^{n+1}_{n+1}-F^{m}{\omega }^{n+1}_{m,{n+1}}-


F^{uv}{\omega }^{n+1}_{uv,n+1}-F^{jkl}{\omega }^{n+1}_{jkl,n+1}&=0, \\


%


dF_i+F_l{\omega }^l_i+F{\omega }^{n+1}_i+F^l{\omega }^{n+1}_{il}


+F^{lk}{\omega }^{n+1}_{ilk}+F^{jkl}{\omega }_{ijkl}^{n+1}&=0, \\


%


dF_0+F_j{\omega }^j_0-F{\omega }^{n+1}_0+


F^j{\omega }_{0j}^{n+1}+F^{jk}{\omega }_{0jk}^{n+1}


+F^{jkl}{\omega }^{n+1}_{0jkl}&=0.


\end{align*}





Из этих уравнений видно, что указанные коэффициенты 


являются компонентами поля геометрического объекта,


заданного на $J^4E$ . Подобъектами этого объекта служат


$\{F^{jkl},F^{jk},F^j\}$, $\{F^{jkl},F^{jk}\}$, $\{F^{jkl}\}$; при этом


$\{F^{jkl}\}$ --- относительный тензор. Коэффициенты $F^i$ и $F^0$ входят в


фундаментальный объект, поле которого задано на многообразии $J^4E$.


Однако в то же самое время $F^i$ и $F^0$ являются также


компонентами самостоятельного объекта,


заданного на многообразии $\ov{J}^4E$, 


главными формами которого служат


$\omega^{\Hat{i}}$, $\omega^{n+1}$, $\omega^{n+1}_{\Hat{j}}$,


$\omega^{n+1}_{\Hat{j}\Hat{k}}\dotsc$


(Это следует из уравнений для $F^i$ и $F^0$.)\;


$\ov{J}^4E$ является фактормногообразием многообразия


$J^4E$ голономных 4-струй сечений расслоенного многообразия $E$.


Следуя [1],


{\it геометрией} дифференциального уравнения (2) будем называть свойства


дифференциально-геометрической структуры, 


определяемой заданием фундаментального объекта.





Из вышеизложенного следует, что геометрия уравнения (2) 


инвариантна относительно допустимых


преобразований локальных координат вида (3).


В дальнейшем для каждого сечения $\sigma\subset E$,


задаваемого уравнением $u=u(t,x^1,\dotsc,x^n)$, будем


рассматривать его поднятия


$\sigma^r\subset J^rE$, задаваемые уравнениями


$u=u(t,x^1,\dotsc,x^n)$,


$\lambda_{\Hat{i}_1\dotsc\Hat{i}_k}=u_{\Hat{i}_1\dotsc\Hat{i}_k}.$





Поднятия $\sigma^{k+1}\subset J^{k+1}E$ сечений


$\sigma\subset E$ и только они являются 


интегральными многообразиями системы уравнений Пфаффа


$$


\omega^{n+1}=\omega^{n+1}_{\Hat{i}}=\dotsb


=\omega^{n+1}_{\Hat{i}_1\dotsc\Hat{i}_k}=0,


$$


где


$\omega^{n+1},\omega^{n+1}_{\Hat{i}},\dotsc,


\omega^{n+1}_{\Hat{i}_1\dotsc\Hat{i}_k}$ ---


контактные формы.





\begin{defnnonum}


Будем называть сечение $\sigma\subset E$ обобщенным 


решением дифференциального уравнения (2), если на поднятии сечения


$\sigma^4\subset J^4E$ дифференциал $d(\lambda_0-f)$ обращается в нуль.


\end{defnnonum}





Решения (2) в собственном смысле содержатся среди обобщенных решений.


Имеет место





\begin{lem}


Если в качестве главных форм многообразия выбраны контактные формы, то


коэффициенты $F_0$ и $F_i$ обращаются в нуль на поднятиях


$\sigma^4\subset J^4E$ сечений


$\sigma\subset E$ тогда и только тогда,


когда сечения $\sigma\subset E$ являются


обобщенными решениями.


\end{lem}





{\bf Доказательство} аналогично случаю уравнений второго порядка,


разобранному в [1].





\section{Понятие связности, определяющей представление нулевой кривизны}





Вкратце напомним основные идеи, связанные с этим понятием.


Ранее достаточно


подробные построения приводились в [1].





Над многообразием $\ov{J}^3E$ как над базой можно построить 


главное расслоение $H(\ov{J}^3E,G)$ с


$r$-параметрической структурной группой Ли $G.$ Структурные формы $\omega^A$


($A=1,\dotsc,r$) этого расслоения удовлетворяют структурным уравнениям


$$


d\omega^A=\frac12 C^A_{BC}\omega^B\wedge\omega^C+


\omega^{\delta}\wedge\omega_{\delta}^A,


$$


где $\omega^{\delta}$ --- главные формы 


многообразия $\ov{J}^3E.$ В этом равенстве


$C^A_{BC}$ --- структурные константы группы $G;$


они антисимметричны по нижним индексам и


удовлетворяют тождествам Якоби


$$


C^A_{BC}C^B_{LM}+C^A_{BL}C^B_{MK}+C^A_{BM}C^B_{KL}=0.


$$





Рассмотрим связность в главном 


расслоении $H(\ov{J}^3E,G),$ определенную заданием поля


объекта связности с компонентами $\Gamma^A_{\delta},$


которые удовлетворяют уравнениям


$d\Gamma^A_\delta+C^A_{BC}\Gamma^B_\delta\omega^C-


\Gamma^A_\varepsilon\omega^\varepsilon_\delta-\omega^A_\delta=


\Gamma^A_{\delta\varepsilon}\omega^\varepsilon$,


где формы $\omega^\varepsilon_\delta$ определяются из уравнений


$d\omega^\varepsilon=\omega^\delta\wedge\omega^\varepsilon_\delta.$


Формы связности имеют вид


$\wt{\omega}^A=\omega^A+\Gamma^A_\delta\omega^\delta$.


Они удовлетворяют структурным уравнениям


$d\wt{\omega}^A=\frac12 C^A_{BC}\wt{\omega}^B\wedge\wt{\omega}^C+\Omega^A$,


где $\Omega^A=R^A_{\delta\varepsilon}\omega^\delta\wedge\omega^\varepsilon$


--- формы кривизны.


Компоненты тензора кривизны имеют вид


$R^A_{\delta\varepsilon}=-\frac12


(\Gamma^A_{[\delta\varepsilon]}+C^A_{BC}\Gamma^B_\delta\Gamma^C_\varepsilon)$.





Можно доказать, что компоненты тензора кривизны вида


$R^A_{\Hat{k}l}$ являются компонентами


тензорного поля, заданного на многообразии $\ov{J}^4E.$


Поэтому обращение компонент


$R_{\Hat{k}l}^A$ в нуль на многообразии $\ov{J}^4E$


либо на каком-то его подмногообразии


имеет инвариантный характер.





Будем говорить, что


{\it связность определяет представление нулевой кривизны для уравнения}


(2), если формы кривизны $\Omega^A$ обращаются в нуль на поднятиях сечений


$\sigma\subset E$ тогда и только тогда, когда сечения $\sigma\subset E$


являются (обобщенными) решениями.





Справедлива





\begin{lem}


Для того чтобы связность, заданная в


расслоении $H(\ov{J}^3E,G),$ определяла


представление нулевой кривизны, необходимо и


достаточно, чтобы компоненты $R^A_{kl}$ и $R^A_{0k}$


обращались в нуль на поднятиях


сечений $\sigma\subset E$ тогда и только тогда,


когда сечения $\sigma\subset E$ являются


$($обобщенными$)$ решениями дифференциального уравнения $(2)$.


\end{lem}





{\bf Доказательство} аналогично случаю уравнения второго


порядка, приведенному в [1].





\section{Проблема существования связности,\protect\\


определяющей представление нулевой кривизны}





Pассмотрим $(N+2)$-мерное $(N>n)$ расслоенное 


многообразие $\cal E$


с расслоенной $(N+1)$-мерной базой $\cal M$ и проекцией $\pi$ (конструкция


аналогична построенной для уравнения (2) в \S\,1).


Пусть


$\vartheta^{\Hat I}$, $\vartheta^{N+1}$,


$\vartheta^0_0$, $\vartheta^I_{\Hat{J}}$,


$\vartheta^{N+1}_{N+1}$, $\vartheta^{N+1}_{\Hat J}$,


$\vartheta^0_{00},\dotsc$


--- структурные формы расслоения голономных реперов 


многообразия $\cal E$


$(I,J=1,\dotsc ,N$; ${\wh I},{\wh J}=0,1,\dotsc ,N).$


Будем предполагать, что на


многообразии $\ov{J}^3E$ задано поле объекта связности с компонентами


$\Gamma^I_{JK}$, $\Gamma^I_{J0}$, $\Gamma^I_{00}$, $\Gamma^0_{00}$,


удовлетворяющими дифференциальным уравнениям


\begin{equation}


\begin{cases}


d\Gamma^I_{JK}+\Gamma^M_{JK}\ov{\vartheta}^I_M-


\Gamma^I_{MK}\ov{\vartheta}^K_J-\Gamma^I_{JM}\ov{\vartheta}^M_K-


\ov{\vartheta}^I_{JK}=0,\\


%


d\Gamma^I_{J0}+\Gamma^M_{J0}\ov{\vartheta}^I_M-


\Gamma^I_{M0}\ov{\vartheta}^M_J-\Gamma^I_{J0}\ov{\vartheta}^0_0-


\Gamma^I_{JM}\ov{\vartheta}^M_0-\ov{\vartheta}^I_{J0}=0, \\


%


d\Gamma^0_{00}-\Gamma^0_{00}\ov{\vartheta}^0_0-\ov{\vartheta}^0_{00}=0, \\


%


d\Gamma^I_{00}+\Gamma^0_{00}\ov{\vartheta}^I_0+


\Gamma^M_{00}\ov\vartheta^I_M-2\Gamma^I_{00}\ov{\vartheta}^0_0-


2\Gamma^I_{M0}\ov{\vartheta}^M_0-\ov{\vartheta}^I_{00}=0.


\end{cases}


\end{equation}


Формы связности имеют вид


\begin{align}


\theta^I_J&=\vartheta^I_J+\Gamma^I_{J0}\vartheta^0+


\Gamma^I_{JK}\vartheta^K, \notag \\


%


\theta^I_0&=\vartheta^I_0+\Gamma^I_{00}\vartheta^0+


\Gamma^I_{J0}\vartheta^J, \\


%


\theta^0_0&= \vartheta^0_0+\Gamma^0_{00}\vartheta^0. \notag


\end{align}





Построение подобного объекта связности для 


эволюционных уравнений второго порядка проведено в


[2], для уравнений третьего порядка --- в [3].





Из уравнений (5) следует, что объект связности содержит 


подобъект с компонентами


$\Gamma^I_{JK}$, $\Gamma^I_{J0}.$


Этому подобъекту сопоставляется связность с формами


$\theta^I_J$, которые удовлетворяют структурным уравнениям


$$


d\theta^I_J=\theta^M_J\wedge\theta^J_M+\Omega^I_J,


$$


где $\Omega^I_J$ --- формы кривизны, имеющие вид


\begin{multline*}


\Omega^I_J=R^I_{JKL}\vartheta^K\wedge\vartheta^L+2R^I_{J0K}


\vartheta^0\wedge\vartheta^K+


2R^I_{J\Hat{K},N+1}\vartheta^{\Hat{K}}\wedge\vartheta^{N+1}+ \\


%


+ 2R^{I\Hat{L}}_{J\Hat{K}}\vartheta^{\Hat


K}\wedge\vartheta_{\Hat{L}}^{N+1}+2R^{ILM}_{J\Hat{K}}


\vartheta^{\Hat{K}}\wedge\vartheta^{N+1}_{LM}+


2R^{ILMN}_{J\Hat{K}}\vartheta^{\Hat{K}}\wedge\vartheta^{N+1}_{LMN}.


\end{multline*}





В многообразии $\cal E$ рассмотрим $(n+2)$-мерное подмногообразие


$\Sigma,$ удовлетворяющее условиям


\begin{itemize}


\item[(а)] одномерный слой расслоения $\cal E$,


проходящий через текущую точку


подмногообразия $\Sigma,$ содержится в $\Sigma$;


\item[(б)] проекция $\pi\Sigma$ подмногообразия


$\Sigma\subset \cal E$ является


$(n+1)$-мерным подмногообразием базы $\cal M$.


\end{itemize}





В этом случае дифференциальные уравнения 


подмногообразия $\Sigma$ могут быть записаны в


виде


\begin{equation}


\vartheta^\alpha=0 \quad (\alpha=n+1,\dotsc,N).


\end{equation}





Продолжим эти уравнения. В результате будут получены соотношения вида


\begin{equation}


\vartheta^\alpha_{\Hat{j}}=\lambda^\alpha_{\Hat{j}\Hat{k}}


\vartheta^{\Hat{k}}.


\end{equation}


Коэффициенты $\lambda^\alpha_{\Hat{j}\Hat{k}}$ 


симметричны по нижним индексам. Зададим


систему координат так, чтобы


$\Gamma^\alpha_{\Hat{j}\Hat{k}}$ обращались в нуль на


$\Sigma\subset \cal E$.


Рассмотрим суммы


$\lambda^\alpha_{\Hat{j}\Hat{k}}+\Gamma^\alpha_{\Hat{j}\Hat{k}}.$


Для них составим


дифференциальные уравнения вида (5).


Из этих уравнений следует, что на подмногообразии


$\Sigma\subset\cal E$


эти суммы являются компонентами тензорного поля, а


поэтому можно утверждать, что класс подмногообразий


$\Sigma\subset\cal E$, на которых этот тензор обращается в


нуль, выделяется инвариантым образом


\begin{equation}


\lambda^\alpha_{\Hat{j}\Hat{k}}+\Gamma^\alpha_{\Hat{j}\Hat{k}}=0.


\end{equation}


Далее будут рассматриваться только такие $\Sigma,$ на которых выполнено это


условие. На этих подмногообразиях из равенства


$\Gamma^\alpha_{\Hat{j}\Hat{k}}=0$


следует, что $\lambda^\alpha_{\Hat{j}\Hat{k}}=0,$


поэтому $\theta^\alpha_{\Hat{j}}=\vartheta^\alpha_{\Hat{j}}=0.$ Уравнения


$\vartheta^\alpha_{\Hat{j}}=0$ также можно продолжить, в результате


чего придем к соотношениям


$$


\vartheta^\alpha_{\Hat{j}\Hat{k}}=


\lambda^\alpha_{\Hat{j}\Hat{k}\Hat{l}}\vartheta^{\Hat{l}}.


$$


Коэффициенты $\lambda^\alpha_{\Hat{j}\Hat{k}\Hat{l}}$


симметричны по нижним индексам.


Процедуру продолжения можно проделать еще раз, в результате будут


получены уравнения


$$


\vartheta^\alpha_{\Hat{j}\Hat{k}\Hat{l}}=


\lambda^\alpha_{\Hat{j}\Hat{k}\Hat{l}\Hat{m}}\vartheta^{\Hat{m}}.


$$





Подмногообразие $\Sigma\subset\cal E$


можно снабдить оснащением специального типа, задаваемым


уравнениями


\begin{equation}


\vartheta^i_\alpha=\vartheta^{N+1}_\alpha=


\vartheta^{N+1}_{\Hat{j}\alpha}=\vartheta^{N+1}_{\alpha\beta}=


\vartheta^{N+1}_{\alpha\beta\gamma}=0(\Mod \vartheta^\delta_\Sigma),


\end{equation}


где $\vartheta^\delta_\Sigma$ ---


главные формы многообразия $\ov{J}^3\Sigma.$ Построение


оснащения проводится по схеме из [1].


В силу (10) и (7) структурные уравнения


для форм связности $\theta^i_j$ имеют вид


\begin{multline*}


d\theta^i_j=\theta^m_j\wedge\theta^i_m+r^i_{jkl}\vartheta^k


\wedge\vartheta^l+2r^i_{j\Hat{k},N+1}


\vartheta^{\Hat{k}}\wedge\vartheta^{N+1}+2r^{i\Hat{l}}_{j\Hat{k}}


\vartheta^{\Hat{k}}\wedge\vartheta^{N+1}_{\Hat{l}}+ \\


%


+2r^{ilm}_{j\Hat{k}}\vartheta^{\Hat{k}}\wedge\vartheta^{N+1}_{lm}+


2r^{ilmn}_{j\Hat{k}}\vartheta^{\Hat{k}}\wedge\vartheta^{N+1}_{lmn}.


\end{multline*}





Таким образом, доказана





\begin{lem}


Если подмногообразие $\Sigma\subset\cal E$


удовлетворяет условиям {\em (а), (б)}


и $(9)$ и снабжено оснащением, задаваемым уравнениями $(10)$,


то задана индуцированная


этим оснащением связность в главном расслоенном 


пространстве $H(\ov{J}^3\Sigma,Gl(n))$


c формами связности $\theta^i_j.$


\end{lem}





Пусть $E$ -- расслоение общего типа над базой $M$ (см. \S\,1).


Установим диффеоморфизм между расслоениями струй над 


многообразиями $E$ и $\Sigma,$ при котором совпадают


соответствующие главные формы многообразий $\ov{J}^4E$ и $\ov{J}^4\Sigma$:


\begin{gather}


\omega^{\Hat{i}}=\vartheta^{\Hat{i}}, \ \; \omega^{n+1}=\vartheta^{N+1}, \ \;


\omega^{n+1}_{\Hat{j}}=\vartheta^{N+1}_{\Hat{j}},\ \;


\omega^{n+1}_{\Hat{j}\Hat{k}}=\vartheta^{N+1}_{\Hat{j}\Hat{k}}, \\


%


\omega^{n+1}_{0jk}=\vartheta^{N+1}_{0jk}+


\lambda^{\alpha}_{0jk}\vartheta^{N+1}_\alpha, \ \;


\omega^{n+1}_{jkl}=\vartheta^{N+1}_{jkl}+


\lambda^\alpha_{jkl}\vartheta^{N+1}_\alpha, \ \;


\omega^{n+1}_{jklm}=\vartheta^{N+1}_{jklm}+


\lambda^\alpha_{jklm}\vartheta^{N+1}_\alpha.\notag 


\end{gather}


Эти уравнения могут также быть продолжены. 


В результате будут получены новые уравнения,


среди которых, в частности, будут следующие:


\begin{equation}


\omega^0_0=\vartheta^0_0, \ \;


\omega^i_0=\vartheta^i_0, \ \; \omega^i_j=\vartheta^i_j.


\end{equation}





В силу (11) и (12) дифференциальные уравнения для 


коэффициентов $F^0$ и $F^i,$ которые содержатся среди компонент 


фундаментального объекта, определяющего геометрию уравнения


(2), становятся уравнениями компонент объекта, 


определенного на многообразии $\ov{J}^4\Sigma.$


Поэтому при установлении диффеоморфизма между расслоениями струй над


$E$ и $\Sigma$ происходит перенос объекта с компонентами $F^0$ и $F^i$


с одного многообразия на другое. Аналогичное замечание справедливо 


относительно произвольного


тензора с компонентами $T^i_{jk}$, $T^I_{j0},$ 


поле которого задано на многообразии


$\ov{J^4}\Sigma,$ а компоненты удовлетворяют уравнениям


$$


\begin{cases}


dT^i_{jk}+T^m_{jk}\ov{\omega}^i_m-T^i_{mk}\ov{\omega}^m_j-


T^i_{jm}\ov{\omega}^m_k=0,\\


%


dT^i_{j0}+T^m_{j0}\ov{\omega}^i_m-T^i_{m0}\ov{\omega}^m_j-


T^i_{j0}\ov{\omega}^0_0-T^i_{jm}\ov{\omega}^m_0=0,


\end{cases}


$$


а также относительно тензора с компонентами


\begin{equation}


r^i_{jkl}=T^i_{j[k}F_{l]}, \quad r^i_{j0k}=T^i_{j[0}F_{k]}.


\end{equation}





Возникает задача: задать оснащение подмногообразия


$\Sigma\subset\cal E$ уравнениями


$$


\vartheta^i_\alpha=\vartheta^{N+1}_\alpha=


\vartheta^{N+1}_{\Hat{j}\alpha}=\vartheta^{N+1}_{\alpha\beta}=


\vartheta^{N+1}_{\alpha\beta\gamma}=0(\Mod\vartheta^\delta_\Sigma)


$$


так, чтобы разложения


для $d\theta^i_j$ имели бы вид (с учетом (13))


\begin{multline*}


d\theta^i_j=\theta^m_j\wedge\theta^i_m+


r^i_{jkl}\omega^k\wedge\omega^l+2r^i_{j0k}\omega^0\wedge\omega^k+


2y^i_{j\Hat{k},n+1}\omega^{\Hat{k}}\wedge\omega^{n+1}+


2y^{i\Hat{l}}_{j\Hat{k}}\omega^{\Hat{k}}\wedge\omega^{n+1}_{\Hat{l}}+ \\


%


+2y^{ilm}_{j\Hat{k}}\omega^{\Hat{k}}\wedge\omega^{n+1}_{lm}+


2y^{ilmn}_{j\Hat{k}}\omega^{\Hat{k}}\wedge\omega^{n+1}_{lmn}.


\end{multline*}


Составим систему дифференциальных уравнений задачи. 


Для этого с учетом (7), (11) и (12)


выпишем структурные уравнения для форм $\theta^i_j,$


содержащихся среди форм связности


(6), после чего вычтем их из соответствующих 


уравнений (13). Таким образом, получится


система дифференциальных уравнений вида


\begin{multline}


R^{i\alpha}_{j\Hat{k}}\omega^{\Hat{k}}\wedge\vartheta^{N+1}_\alpha+


R^{il\alpha}_{j\Hat{k}}\omega^{\Hat{k}}\wedge\vartheta^{N+1}_{\alpha\beta}+


R^{i\alpha\beta\gamma}_{j\Hat{k}}\omega^{\Hat{k}}\wedge


\vartheta^{N+1}_{\alpha\beta\gamma}=\frac{1}{2}\rho^i_{jkl}


\omega^k\wedge\omega^l+\rho^i_{j0k}\omega^0\wedge\omega^k+ \\


%


+\chi^i_{j\Hat{k},n+1}\omega^{\Hat{k}}\wedge\omega^{n+1}+


\chi^{i\Hat{l}}_{j\Hat{k}}\omega^{\Hat{k}}\wedge\omega^{n+1}_{\Hat{l}}+


\chi^{ilm}_{j\Hat{k}}\omega^{\Hat{k}}\wedge\omega^{n+1}_{lm}+


\chi^{ilmn}_{j\Hat{k}}\omega^{\Hat{k}}\wedge\omega^{n+1}_{lmn},


\end{multline}


где


\begin{alignat*}{2}


\rho^i_{jkl}&=r^i_{jkl}-R^i_{jkl}, &\quad\rho^i_{j0k}&=r^i_{j0k}-R^i_{j0k},\\


%


\chi^i_{j\Hat{k},n+1}&=y^i_{j\Hat{k},n+1}-R^i_{j\Hat{k},n+1}, & \quad


\chi^{i\Hat{l}}_{j\Hat{k}}&=y^{i\Hat{l}}_{j\Hat{k}}-R^{i\Hat{l}}_{j\Hat{k}}, \\


%


\chi^{ilm}_{j\Hat{k}}&=y^{ilm}_{j\Hat{k}}-R^{ilm}_{j\Hat{k}}, &\quad


\chi^{ilmn}_{j\Hat{k}}&=y^{ilmn}_{j\Hat{k}}-R^{ilmn}_{j\Hat{k}}.


\end{alignat*}





При этом коэффициенты $\rho^i_{jkl}$, $\rho^i_{j0k}$ заданы заранее, а


$\chi^i_{j\Hat{k},n+1}$, $\chi^{il}_{j\Hat{k}}$, $\chi^{ilm}_{j\Hat{k}}$,


$\chi^{ilmn}_{j\Hat{k}}$ --- неизвестные функции.


Если будет доказана совместность системы (14),


то таким образом будет доказано


существование представления нулевой кривизны для уравнения (2).


Подвергая уравнения (14)


внешнему дифференцированию и присоединяя полученные уравнения


к исходным, получаем замкнутую систему дифференциальных уравнений задачи.





Для исследования совместности системы используется метод 


Кэлера ([4], с.\,207). В случае $n=1$


выясняется, что совместность имеет место, если компоненты объекта связности


$\Gamma^I_{JK}$, $\Gamma^I_{J0}$ выбраны так,


что тензор кривизны связности на


многообразии $\Sigma\subset\cal E$ удовлетворяют условиям


\begin{itemize}


\item[a)] можно указать наборы чисел: $\alpha_1<\alpha_2<\alpha_3<\alpha_4,$


$\beta_1<\beta_2<\beta_3<\beta_4,$


$\gamma_1<\gamma_2<\gamma_3<\gamma_4$ такие, что


$$


\begin{vmatrix}


R^{1\beta_1\alpha_1\gamma_1}_{10} &


R^{1\beta_2\alpha_1\gamma_1}_{10} &


R^{1\beta_3\alpha_1\gamma_1}_{10} &


R^{1\beta_4\alpha_1\gamma_1}_{10} \\


\\


R^{1\beta_1\alpha_2\gamma_2}_{10} &


R^{1\beta_2\alpha_2\gamma_2}_{10} &


R^{1\beta_3\alpha_2\gamma_2}_{10} &


R^{1\beta_4\alpha_2\gamma_2}_{10} \\


%


\dotsb & \dotsb & \dotsb & \dotsb \\


%


R^{1\beta_1\alpha_4\gamma_4}_{10} &


R^{1\beta_2\alpha_4\gamma_4}_{10} &


R^{1\beta_3\alpha_4\gamma_4}_{10} &


R^{1\beta_4\alpha_4\gamma_4}_{10}


\end{vmatrix}


\ne 0;


$$


\item[b)] можно указать компоненты $R^{1\lambda_1\mu_1\gamma_1}_{10}$,


$R^{1\lambda_2\mu_2\gamma_2}_{10}$, $R^{1\lambda_1\mu_1\gamma_1}_{11}$,


$R^{1\lambda_2\mu_2\gamma_2}_{11},$ удовлетворяющие условию


$$


\begin{vmatrix}


R^{1\lambda_1\mu_1\gamma_1}_{10} & R^{1\lambda_2\mu_2\gamma_2}_{10} \\


\\


R^{1\lambda_1\mu_1\gamma_1}_{11} & R^{1\lambda_2\mu_2\gamma_2}_{11} \\


\end{vmatrix}


\ne 0;


$$


\item[c)] можно указать числа


$\lambda_i$, $\mu_i$, $\gamma_k$, $\theta_l$ $(i,j,k,l=1,\dotsc,6$)


такие, что выполняется


$$


\begin{vmatrix}


R^{\gamma_1\lambda_1\mu_1\theta_1}_{11} &


\dotsc & R^{\gamma_6\lambda_1\mu_1\theta_1}_{11} \\


. & \dotsc & . \\


. & \dotsc & . \\


. & \dotsc & . \\


R^{\gamma_1\lambda_6\mu_6\theta_6}_{11} &


\dotsc & R^{\gamma_6\lambda_6\mu_6\theta_6}_{11}


\end{vmatrix}


\ne 0.


$$


\end{itemize}





Таким образом, имеет место





\begin{thmnonum}


Пусть


\begin{itemize}


\item[1)] задано эволюционное дифференциальное уравнение 


третьего порядка $(2);$


\item[2)] на многообразии $\ov{J}^3\varepsilon$, где $\varepsilon$ ---


$(N+2)$-мерное $(N>n)$ расслоенное многообразие 


с $(N+1)$-мерной расслоенной базой,


задано поле объекта связности с компонентами


$\Gamma^I_{JK}$, $\Gamma^I_{J0}$,


$\Gamma^I_{00}$, $\Gamma^0_{00}$, удовлетворяющими


дифференциальным уравнениям $(6);$


\item[3)] в многообразии $\varepsilon$ задано $(n+2)$-мерное подмногообразие


$\Sigma\subset\varepsilon$, удовлетворяющее условиям {\em (а), (б)} и


$(8);$


\item[4)] между расслоениями струй над 


многообразиями $\Sigma$ и $\varepsilon$


установлен диффеоморфизм, при котором совпадают главные формы в


$\ov{J}^4\varepsilon$ и $\ov{J}^4\Sigma$.


\end{itemize}





Тогда в случае $n=1$ при выполнении условий {\em a), b), c)}


можно найти такое оснащение подмногообразия $\Sigma,$ что


индуцируемая им связность будет определять представление


нулевой кривизны для дифференциального уравнения $(2)$.


\end{thmnonum}





\begin{notenonum}


Связность с компонентами $\Gamma^I_{JK}$, $\Gamma^I_{J0}$,


$\Gamma^I_{00}$, $\Gamma^0_{00}$ выбирается произвольно,


поэтому ее всегда можно выбрать


так, чтобы выполнялись условия a), b), c). Следовательно, представление


нулевой кривизны существует для любого эволюционного уравнения третьего


порядка с одной пространственной переменной.


\end{notenonum}





Данное утверждение может быть использовано при решении вопроса о существовании


законов сохранения для дифференциального уравнения третьего порядка.





\section{Псевдопотенциалы и законы сохранения}





Впервые понятие псевдопотенциала появилось в работах


Ф.\,Эстабрука и Х.\,Уолквиста [5] и [6].


Подробный анализ связи псевдопотенциалов и законов 


сохранения со связностями, определяющими


представление нулевой кривизны, содержится в [1]. Напомним основные идеи.


Пусть задано эволюционное дифференциальное уравнение


третьего порядка.


Законы сохранения, ассоциированные с уравнением, 


задаются системой 1-форм $\theta^I$ вида


$$


\theta^I=dX^I+F^Idt+G^Idx, \quad I,J=1,\dotsc,N,


$$


где $X^I$ --- некоторые переменные, а $x$


--- единственная пространственная переменная. 


Функции $F^I$ и $G^I$ зависят от


переменных $X^1,\dotsc,X^N$, $t$, $x$, 


неизвестной функции $u$ и частных производных


функции $u$ по $x$. В работе [4]


из соображений простоты построений рассматривались функции $F^I$ и


$G^I$, которые не зависят от $t$ и $x$; при этом


$$


\theta^I=dX^I+F(u,\lambda_1,\lambda_{11},X^I)dt+


G^I(u,\lambda_1,\lambda_{11},X^I)dx.


$$


Если система дифференциальных форм $\theta^I$ 


вполне интегрируема, тогда на каждом


решении $X^I=X^I(x,t)$ системы $\theta^I|_\sigma=0$ функции $F^I$ и


$G^I$ удовлетворяют уравнению


\begin{equation}


\frac{\partial G^I}{\partial t}-\frac{\partial F^I}{\partial x}=0.


\end{equation}


Если известны $N$ форм $\theta^I$ ($N$ законов сохранения в смысле Уолквиста


и Эстабрука), то можно найти $N$ уравнений такого вида.


В том случае, когда $F^I$ и $G^I$ зависят от переменных


$X^I$ линейно, уравнения (15)


принимают классический вид законов сохранения. 


Уолквист и Эстабрук называли переменные $X^I$ {\it


псевдопотенциалами.}





Следуя [1], можно доказать





\begin{propnonum}


Если в главном расслоении $H({\ov J}^{3}E,G)$ задана связность, определяющая


представление нулевой кривизны, ассоциированное с


дифференциальным уравнением $(2)$, и


выбрано сечение, то любому присоединенному 


расслоению $F(H({\ov J}^{3}E,G))$ соответствует серия законов


сохранения.


\end{propnonum}





\begin{notenonum}


Принимая во внимание утверждение теоремы 


предыдущего параграфа, можно сделать вывод,


что для любого эволюционного уравнения третьего порядка с одной


пространственной переменной можно построить 


бесконечно много законов сохранения.


\end{notenonum}





Есть все основания предполагать, что метод, 


использованный впервые в [1] при доказательстве 


теоремы существования связности,


определяющей представление нулевой кривизны 


для эволюционного уравнения второго порядка,


примененный в данной статье к эволюционному 


уравнению третьего порядка, может быть


успешно использован при рассмотрении уравнений в частных производных


произвольного типа более высоких порядков.





\begin{thebibliography}{9}





\bibitem{1}


Рыбников А.К. {\em О геометрии псевдопотенциалов эволюционных 


уравнений} // Изв.


вузов. Математика. -- 1995. -- \No\,5. -- C.\,55--67.


\bibitem{2}


Рыбников А.К. {\em О геометрии эволюционного уравнения второго порядка} //


Вестн. Моск. ун-та. Cер. матем., механ. --


1994. -- \No\,2. -- C.\,79--85.


\bibitem{3}


Семенов К.В. {\em О геометрии эволюционного уравнения третьего порядка} //


Изв. вузов. Математика. -- 1998. -- \No\,6. -- C.\,67--74.


\bibitem{4}


Фиников С.П. {\em Метод внешних форм Картана в дифференциальной


геометрии}. -- М.--Л.: ГИТТЛ, 1948. -- 432 с.


\bibitem{5}


Wahlquist H.D., Estabrook F.B. {\em Prolongation structures of nonlinear


evolution equations} // J. Math. Phys. -- 1975. -- V.\,16.


-- \No\,1. -- P.\,1--7.


\bibitem{6}


Estabrook F.B., Wahlquist H.D. {\em Prolongation structures of nonlinear


evolution equations} // J. Math. Phys. -- 1976. --


V.\,17. -- \No\,7. -- P.\,1293--1297.





\end{thebibliography}





\vskip 0.5cm





\post{\it Московский государственный университет}{\it Поступила}


\post{}{10.09.1999}





\label{end}


\end{document}


